\section{Cramer's Rule}
\label{sec:cramers-rule}

\noindent\textbf{Subject:} Linear Algebra for ML (3Blue1Brown) \\
\textbf{Status:} growing

\subsection*{Summary}
Cramer's Rule is a way to solve linear equations geometrically using determinants.

\subsection{Core Idea}

\begin{defbox}
It's a way to \textbf{solve linear equations} geometrically.
\end{defbox}

\[
A \cdot x = B
\]

So geometrically it would look like:

\[
\begin{bmatrix} 3 & 2 \\ 4 & 2 \end{bmatrix} \begin{bmatrix} x \\ y \end{bmatrix} = \begin{bmatrix} -4 \\ -2 \end{bmatrix}
\]

\begin{itemize}[nosep]
  \item $A$ = transformation matrix
  \item Mystery input vector $\rightarrow$ transformed vector in our language
\end{itemize}

If $T(\vec{v}) \cdot T(\vec{w}) = \vec{v} \cdot \vec{w}$ for all $\vec{v}$ and $\vec{w}$:

\begin{defbox}
$T$ is \textbf{``Orthonormal''} --- it leaves all the basis vectors perpendicular with unit lengths.
\end{defbox}

\subsection*{Related Topics}
\begin{itemize}[nosep]
  \item Section~\ref{sec:cross-product} --- Cross Product
  \item Section~\ref{sec:change-of-basis} --- Change of Basis
\end{itemize}

\bigskip
\noindent\textit{Tags: linear-algebra, cramers-rule, math, phase-1}
